\section{Hilbert-adjoint operators, symmetric and self-adjoint linear operators}


\begin{exercise}{1}
Show that a self-adjoint linear operator is symmetric.
\end{exercise}
\begin{proof}
By Lemma 10.2-4, we have that a densely defined linear operator $T$ is symmetric if and only if $T \subseteq T^\ast$.
By definition, $T: \cD(T) \to H$ densely defined in $H$ is self-adjoint if $T = T^\ast$, so that if $T$ is self-adjoint, then it must be symmetric.
\end{proof} 

\begin{exercise}{2}
If $S$ and $T$ are such that $ST$ is densely defined in $H$, show that $T^\ast S^\ast \subseteq (ST)^\ast$, and if $S$ is defined on all of $H$ and is bounded, then $(S T)^\ast = T^\ast S^\ast $
\end{exercise}
\begin{proof}
Let $y \in \cD(S^\ast)$, so that $\brackets{Sx, y} = \brackets{x, y^\ast}$, for all $x \in H$ and $y^\ast = S^\ast y$.
Moreover, if $z \in \cD(T^\ast)$, then $\brackets{Tx, z} = \brackets{x, z^\ast}$, for all $x \in H$ and $z^\ast = T^\ast z$.
Thus, if $y \in \cD(T^\ast S^\ast)$, we have
\begin{align*}
    \brackets{STx, y}
    = \brackets{Tx, S^\ast y}
    = \brackets{x, T^\ast S^\ast y},
\end{align*}
for all $x \in H$, so that $y \in \cD((ST)^\ast)$.

For the second part, suppose $S$ is defined on all of $H$, so that $S^\ast$ corresponds to the usual adjoint of $S$.
Let $y \in \cD((ST)^\ast)$, so that $\brackets{ST x, y} = \brackets{Tx, S^\ast y} = \brackets{x, y^\ast}$, for all $x \in H$ and $y^\ast = T^\ast (S^\ast y)$.
But note again, this is precisely the same as $y \in \cD(T^\ast S^\ast)$, completing the proof.
\end{proof} 

\begin{exercise}{3}
Let $H$ be a complex Hilbert space and $T: \cD(T) \to H$ linear and densely defined in $H$.
Show that $T$ is symmetric if and only if $\brackets{Tx, x}$ is real for all $x \in \cD(T)$.
\end{exercise}
\begin{proof}
We will follow the same strategy as in 3.10-3.

($\Rightarrow$)
Suppose $T$ is symmetric, so that for all $x, y \in \cD(T)$, we have $\overline{\brackets{Tx, y}} = \brackets{y, Tx} = \brackets{Tx, y}$. 
In particular, this holds for $y = x$, so that $\overline{\brackets{Tx, x}} = \brackets{Tx, x}$;
that is, $\brackets{Tx, x}$ is real for all $x$.

($\Leftarrow$)
Suppose $\brackets{Tx, x}$ is real for all $x \in \cD(T)$, then 
\begin{align*}
    \brackets{Tx, x} 
    = \overline{\brackets{Tx, x}} 
    = \overline{\brackets{x, T^\ast x}}
    = \brackets{T^\ast x, x}.
\end{align*}
Thus $\brackets{(T - T^\ast)x, x} = 0$, so that $T = T^\ast$ by Lemma 3.9-3.b  (since $H$ is complex), proving the desired result.

\end{proof} 

\begin{exercise}{4}
If $T$ is symmetric, show that $T^{\ast \ast}$ is symmetric.
\end{exercise}
\begin{proof}
Since $T$ is symmetric, we know $T \subseteq T^\ast$ (10.2-4), so that $T^{\ast \ast} \subseteq T^\ast$ (10.2-1.a), and also $T^{\ast \ast} \subseteq T^{\ast \ast \ast}$, applying the same rule again.
Letting $T^{\ast \ast} = S$, we have $S \subseteq S^\ast$, so that by 10.2-4 we get that $S = T^{\ast \ast}$ is symmetric too.
\end{proof} 

\begin{exercise}{5}
If a linear operator $T$ is densely defined in $H$ and its adjoint is defined on all of $H$, show that $T$ is bounded.
\end{exercise}
\begin{proof}
We have that $T \subseteq T^{\ast \ast}$ (10.2-1.b), and we proved in exercise 10.1.9 that if an operator is defined everywhere, then its adjoint is bounded, so that $T^{\ast \ast}$ is bounded.
Since $T^{\ast \ast}$ is an extension of $T$, then it must be the case that $T$ itself is bounded, completing the proof.
\end{proof} 

\begin{exercise}{6}
Show that $y = (y_n) = Tx = (x_n / n)$ defines a bounded self-adjoint linear operator $T: \ell^2 \to \ell^2$ which has an unbounded self-adjoint inverse.
\end{exercise}
\begin{proof}
We will first prove that $T$ is bounded.
Note,
\begin{align*}
    \norm{Tx}^2
    = \sum \frac{\absoluteValue{x_n}^2}{n^2}
    \leq \sum \absoluteValue{x_n}^2
    = \norm{x}^2,
\end{align*}
so that $\norm{T} \leq 1$.
Now consider $x = (1, 0, \dots)$, so that $\norm{Tx} = 1$, giving us that $\norm{T} = 1$.
Furthermore, 
\begin{align*}
    \brackets{Tx, y} 
    = \sum (x_n / n) \bar{y_n} 
    = \sum x_n (\bar{y_n} / n)
    = \brackets{x, Ty}, 
\end{align*}
so that $T$ is self-adjoint.

The inverse of $T$ is $T^{-1}y = (x_n)$, where $x_n = n y_n$, which we proved is unbounded in exercise 10.1.8.
\end{proof} 

\begin{exercise}{7}
Let $T: \cD(T) \to H$ be a bounded symmetric linear operator.
Show that $T$ has a bounded symmetric linear extension $\bar{T}$ to $\overline{\cD(T)}$.
\end{exercise}
\begin{proof}
Since $T$ is bounded we can extend it to a bounded linear operator $\bar{T}: \overline{\cD(T)} \to H$ (2.7-11).
To see $\bar{T}$ is symmetric, let $x, y \in \overline{\cD(T))}$, so that there exists sequences $(x_n)$ and $(y_n)$ in $\cD(T)$ with $x_n \to x$ and $y_n \to y$.
We then have $\brackets{\bar{T}x_n, y_n} = \brackets{x_n, \bar{T}y_n}$ since $T$ is symmetric.
Taking limits as $n \to \infty$ and the fact that the inner product is continuous gives us that $\bar{T}$ is indeed symmetric, as required.
\end{proof} 

\begin{exercise}{8}
If $T$ is symmetric and $\bar{T}$ is a symmetric extension of $T$, show that $\bar{T} \subseteq T^\ast$.
\end{exercise}
\begin{proof}
By hypothesis, $T \subseteq \bar{T}$, so that by 10.2-1, we have $\bar{T}^\ast \subseteq T^\ast$.
Furthermore, since $\bar{T}$ is symmetric, we have that $\bar{T} \subseteq \bar{T}^\ast$ (10.2-4).
Putting these results together, we have that $\bar{T} \subseteq \bar{T}^\ast \subseteq T^\ast$, implying the desired result.
\end{proof} 

\begin{exercise}{9}
A maximally symmetric linear operator is defined to be a symmetric linear operator which has no proper symmetric extensions.
Show that a self-adjoint linear operator $T$ is maximally symmetric.
\end{exercise}
\begin{proof}
This is a corollary to the previous exercise. 
To see this, suppose $T$ is a maximally symmetric operator and $\bar{T}$ is a symmetric extension of $T$.
Thus, $T^\ast = T \subseteq \bar{T}$ and by the previous exercise $\bar{T} \subseteq T^\ast = T$.
Hence, it must be the case that $\bar{T} = T$, as required.
\end{proof} 

\begin{exercise}{10}
If a self-adjoint linear operator $T: \cD(T) \to H$ is injective, show that 
\begin{enumerate}
    \item $\overline{\cR(T)} = H$ and
    \item $T^{-1}$ is self-adjoint.
\end{enumerate}
\end{exercise}
\begin{proof}
\begin{enumerate}
    \item First, note that since $T$ is injective, then $\nullspace(T) = \set{0}$.
    
    Furthermore, $\nullspace(T^\ast) = \cR(T)^\perp$.
    To see this, let $y \in \nullspace(T^\ast)$ with $y \neq 0$, then for all $x \in \cD(T)$, we have $\brackets{Tx, y} = \brackets{x, T^\ast y} = 0$, so that $y \in \cR(T)^\perp$ and $\nullspace(T^\ast) = \cR(T)^\perp$.
    Conversely, let $y \in \cR(T)^\perp$, then for all $x$, we have $0 = \brackets{Tx, y} = \brackets{x, T^\ast y}$, so that $y \in \nullspace(T^\ast)$ and $\cR(T)^\perp \subseteq \nullspace(T^\ast)$.

    Since $T = T^\ast$, given that $T$ is self-adjoint, we have that $\nullspace(T) = \nullspace(T^\ast)$.

    Putting these together, we have that $\cR(T)^\perp = \nullspace(T^\ast) = \nullspace(T) = \set{0}$.
    The closure of $\cR(T)^\perp$ is simply $\set{0}$, meaning that $\overline{\cR(T)} = H$, as required.
    
    \item By 10.2-2, we have that $T^\ast$ is injective and $(T^\ast)^{-1} = (T^{-1})^\ast$, however, since $T = T^\ast$ due to $T$'s self-adjointness, we have $T^{-1} = (T^{-1})^\ast$, which means precisely that $T^{-1}$ is self-adjoint.
\end{enumerate}
\end{proof} 
